\chapter{Stima spettrale}

\begin{figure}[h]
    \centering
    \includegraphics[scale = 0.5]{SpettrogrammaVocaliAeI.jpg}
\end{figure}  

\newpage 

\section{Stima dello spettro}

\footnote{ 
    Slide del prof | 13 - Stima spettrale | pag 2 \\
    Appunti | 2025-12-09 | pag 1 \\  
    Video Lezione CADSP Lunedì 12 Dicembre 2022 | 30:45 - 44:20  
}

\begin{tcolorbox}
    La stima dello spettro è quello che realmente faremo SEMPRE perchè i segnali analogici reali sono infiniti (o di lunghezza molto lunga rispetto alla memoria finita del calcolatore). \newline 
    
    Mi dispiace che ho dovuto scrivere più di 300 pagine per arrivare a questo punto, ma è stato doveroso perchè ci servono delle basi per capire cosa faremo realmente: 
    è come studiare e passare l'esame di Analisi Matematica 1, per poi utilizzare, ogni giorno, massimo le somme, sottrazioni, divisioni, moltiplicazioni e approssimazioni. \newline 
    
    Questo è il primo "vero" capitolo di DSP, su cui si basa anche l'esame finale, specialmente la prova su Matlab. \newline  

    Buona lettura e buono studio. 
\end{tcolorbox}

In molte applicazioni si è interessati al calcolo dello spettro di un segnale campionato, in particolare si vuole fare una stima dello spettro di densità di energia o lo spettro di densità di potenza. \newline 

Si parla di stima dello spettro e non di spettro reale perchè, nella realtà, avendo a che fare con segnali casuali, non possiamo conoscere a priori il segnale, e quindi il massmo che possiamo fare è stimarlo. \newline 

Inoltre si parla di stima perchè le sequenze hanno una durata che può essere: 

\begin{itemize}
    \item finita 
    \item infinita (o molto lunga per le risorse limitate)
\end{itemize}

Andiamo ad affrontare queste due casistiche nel dettaglio. \newline 

\newpage 

\subsection{Sequenze di durata finita}

\footnote{ 
    Slide del prof | 13 - Stima spettrale | pag 2 - 3\\
    Appunti | 2025-12-09 | pag 1 - 2  \\ 
    Video Lezione CADSP Lunedì 12 Dicembre 2022 | 44:20 - 47:23   
}

Consideriamo di avere una sequenza $x \left[ n\right]$ di lunghezza finita, e che quindi ha tutti i valori di $x \left[ n\right]$ per $n = 0, \dots, N-1$. \newline 

Allora, per ogni valore di n, possiamo calcolarci il suo spettro $X \left[ k\right]$ come: 

{
    \Large 
    \begin{equation}
        X \left[ k\right]
        = 
        \left.
            X \left(e^{\jmath \omega}\right)
        \right|_{\omega = \frac{2 \pi}{k} \cdot k}
        \quad 
        \text{ per }
        k = 0, \dots, N - 1
    \end{equation}
}

In questo caso particolare, la FFT fornisce N campioni equispaziati sul cerchio unitario della corrispondente DTFT (senza errori). \newline 

Si noti che, essendo la sequenza $x \left[ n\right]$ finita, non calcoleremo la DTFT, bensì la DFT: 

{
    \Large 
    \begin{equation}
        \begin{split}
            X \left(e^{\jmath \omega}\right) 
            &= 
            \sum_{n = - \infty}^{+ \infty} 
            x \left[ n\right]
            \cdot 
            e^{- \jmath \omega n}
            \\
            &= 
            \sum_{n = 0}^{N + 1} 
            x \left[ n\right]
            \cdot 
            e^{- \jmath \omega n} 
            \quad 
            \text{ per }
            \abs{\omega} < \pi
        \end{split}
    \end{equation}
}

La funzione $X \left(e^{\jmath \omega}\right)$ è una funzione calcolabile analitcamente perchè la sequenza è di durata finita. \newline 

Questo è il caso ideale e non saremmo mai in grado di calcolare analitcamente la funzione $X \left(e^{\jmath \omega}\right)$ nella realtà. \newline 

Come studiato in \ref{ch: FFT} \nameref{ch: FFT} da pag \pageref{ch: FFT}, la FFT è uno strumento molto utile per calcolare la DFT di $x \left[ n\right]$. \newline 

Per avere un campionamento più fitto, nella DFT basta fare uno zero padding, ma nella realtà questo caso non sarà possibile farlo. \newline 

\newpage 

\subsection{Sequenze di durata infinita}

\footnote{ 
    Slide del prof | 13 - Stima spettrale | pag 4 - 8\\
    Appunti | 2025-12-09 | pag 2  \\ 
    Video Lezione CADSP Lunedì 12 Dicembre 2022 | 47:23 - 1:15:40     
}

Questo è il caso di tutte le sequenze molto lunghe o di cui si conosce solo un breve tratto. \newline 

La DTFT di $x \left[ n\right]$, cioè: 

{
    \Large 
    \begin{equation}
        X \left(e^{\jmath \omega}\right)
        = 
        \sum_{n = - \infty}^{+ \infty} 
        x \left[ n\right]
        \cdot 
        e^{-\jmath \omega n}
    \end{equation}
}

non si può calcolare direttamente perchè i valori n-esimi sono infiniti (la sommatoria si svolge nell'intervallo $\left[ - \infty, + \infty\right]$). \newline 

Quindi, l'idea è quella di troncare la sequenza $x \left[ n\right]$ alla lunghezza L, utilizzando una finsitra, proprio come abbiamo studiato in \ref{sec: Progetto di filtri FIR: Windowing} \nameref{sec: Progetto di filtri FIR: Windowing} da pag \pageref{sec: Progetto di filtri FIR: Windowing}, 
ma questa volta non per la risposta impulsiva dei filtri, bensì troncare direttamente la sequenza $x \left[ n\right]$. \newline 

Oltre al troncamente, dobbiamo considerare il caso più semplice, cioè quella di una sequenza $x \left[ n\right]$ stazionaria. \newline 

\begin{tcolorbox}
    La stazionarietà l'avevamo studiata anche nel corso di Teoria dei segnali con il prof. Franco Chiaraluce, ma per un breve ripasso al volo di cosa si intende per stazionarietà: \newline 

    Processo stazionario - Wikipedia\\
    \url{https://it.wikipedia.org/wiki/Processo_stazionario}
\end{tcolorbox}

Il troncamento implica un prodotto tra $x \left[ n\right]$ (di lunghezza infinita) e la finestra $w \left[ n\right]$ (di lunghezza finita). \newline 

Da questo momento in poi, non consideremo più $x \left[ n\right]$ di lunghezza infinita, bensì la sua stima $\hat{x} \left[ n\right]$: 

{
    \Large 
    \begin{equation}
        \hat{x} \left[ n\right] 
        = 
        x \left[ n\right]
        \cdot 
        w \left[ n\right]
    \end{equation}
}

e la finestra $w \left[ n\right]$ deve essere sempre non nullo. \newline 

Calcoliamo $\hat{X} \left(e^{\jmath \omega}\right)$, che è la stima di $X \left(e^{\jmath \omega}\right)$: 

{
    \Large 
    \begin{equation}
        \begin{split}
            \hat{X} \left(e^{\jmath \omega}\right)
            &= 
            \sum_{n = - \infty}^{+ \infty} 
            \hat{x} \left[ n\right]
            \cdot 
            e^{-\jmath \omega n}
            \\
            &= 
            \sum_{n = - \infty}^{+ \infty} 
            x \left[ n\right]
            \cdot
            w \left[ n\right] 
            \cdot
            e^{-\jmath \omega n}
            \\
            &= 
            \sum_{n = 0}^{L - 1} 
            x \left[ n\right]
            \cdot
            w \left[ n\right] 
            \cdot
            e^{-\jmath \omega n}
        \end{split}
    \end{equation}
}

\newpage 

Poichè $\hat{x} \left[ n\right]$ è di durata finita, di lunghezza L (che è la lunghezza della finestra $w \left[ n\right]$), $\hat{X} \left(e^{\jmath \omega}\right)$ non sarà la DTFT, bensì la DFT a N punti: 

{
    \Large 
    \begin{equation}
        \begin{split}
            \hat{X} \left[ k\right]
            &= 
            \left.
            \hat{X} \left(e^{\jmath \omega}\right)
            \right|_{\omega = \frac{2 \pi}{N} \cdot k}
            \\
            &= 
            \sum_{n = 0}^{L - 1} 
            x \left[ n\right]
            \cdot
            w \left[ n\right] 
            \cdot
            e^{-\jmath \frac{2\pi}{N} n k} 
            \quad 
            \text{ per }
            k = 0, \dots, N - 1
        \end{split}
    \end{equation}
}

e quindi $\hat{X} \left(e^{\jmath \omega}\right)$ sarà distorto rispetto a $X \left(e^{\jmath \omega}\right)$ perchè la finestra è più corta di $x \left[ n\right]$, che è lunga N. \newline 

Affinchè la formula di $\hat{X} \left[ k\right]$ sia valida, occorre che: 

{
    \Large 
    \begin{equation}
        N \geq L
    \end{equation}
}

oppure si fa zero padding su $x \left[ n\right]$ se N è minore di L, in modo da arrivare a $N \geq L$. \newline 

Aumentare N non aggiunge l'informazione, non migliora l'informazione, non migliora la stima: permette di distinguere le frequenze. \newline 

Per migliorare la stima, bisognerebbe aumentare L, cioè la lunghezza della finestra perchè abbiamo più informazioni di $x \left[ n\right]$ (o come la chiama il prof "abbiamo più birra"). \newline 

Siccome abbiamo fatto una finestratura, quindi un prodotto nel dominio n, grazie alle proprietà tra tempo e Fourier, 
abbiamo che avremo a che fare con una convoluzione nello spettro, cioè valgono queste relazioni: 

{
    \Large 
    \begin{equation}
        \begin{split}
            \hat{x} \left[ n\right]
            &= 
            x \left[ n\right]
            \cdot 
            w \left[ n\right]
            \\
            &\updownarrow 
            \\
            \hat{X} \left(e^{\jmath \omega}\right)
            &= 
            \dfrac{1}{2 \pi}
            \cdot 
            \int_{- \pi}^{+ \pi}
            X \left(e^{\jmath \theta}\right)
            \cdot 
            W \left(e^{\jmath \left(\omega - \theta\right)}\right)
            d \theta
        \end{split}
    \end{equation}
}

\begin{tcolorbox}[colback=red!5!white,colframe=red!75!black] \label{q: 43}
Siccome nel mondo DSP non esistono integrali, se volessimo fare i precisini,in realtà $\hat{X} \left(e^{\jmath \omega}\right)$ è una stima di una stima di $X \left(e^{\jmath \omega}\right)$ ?
\end{tcolorbox}

Come studiato nel capitolo della progettazione dei filtri, possiamo decidere una tra le seguenti finestre studiate in \ref{subsec: Elenco delle finestre principali} \nameref{subsec: Elenco delle finestre principali} a pag \pageref{subsec: Elenco delle finestre principali}, 
cioè possiamo scegliere una tra le seguenti finestre: 

\begin{itemize}
    \item rettangolare 
    \item Bartlett 
    \item Hamming 
    \item Hanning 
    \item Kaiser 
\end{itemize}

\newpage 

Ma, siccome noi valuteremo una stima di $\hat{X} \left(e^{\jmath \omega}\right)$, perchè non useremo gli integrali bensì una sommatria, allora, il fatto di avere dei valori finiti, 
può dare problemi per la localizzazione di picchi o di punti particolari, come visto in questa figura: 

\begin{figure}[h]
    \centering
    \includegraphics[scale = 0.8]{granulostià esempio.png}
\end{figure}  

Questo problema prende il nome di granulosità, che avevamo già affrontato, implicitamente, con il fenomeno di Gibbs in \ref{subsec: Conseguenze del Windowing} \nameref{subsec: Conseguenze del Windowing} a pag \pageref{subsec: Conseguenze del Windowing}. \newline 

\begin{tcolorbox}[colback=red!5!white,colframe=red!75!black] \label{q: 44}
Spiegazione della slide e con quale forma d'onda è stata fatta la convoluzione la finestra. 
\end{tcolorbox}

Come nela caso della progettazione dei filtri FIR, possiamo considerare i seguenti parametri, detti anche gradi di libertà, per la finestra $w \left[ n\right]$: 

\begin{itemize}
    \item la lunghezza della finestra L, aumentando la lunghezza della finestra, nella DFT avremo un lobo principarle più stretto 
    \item il tipo di finestra, quindi la forma della finestra, perchè attenuazioni dei lobi laterali alta implica un lobo principale più largo 
    \item la lunghezza della DFT (e quindi della FFT) in modo tale che $N > L$, perchè maggiore è N e minore è la granulosità 
\end{itemize}

\newpage 

\subsubsection{Problemi della stima di una sequenza infinita}

\footnote{ 
    Slide del prof | 13 - Stima spettrale | pag 9\\
    Appunti | 2025-12-09 | pag 2 \\ 
    Video Lezione CADSP Lunedì 12 Dicembre 2022 | 1:15:40 - 1:18:45 
}

La stima dello spettro $\abs{X \left(e^{\jmath \omega}\right)}^{2}$ di un segnale $x \left[ n\right]$, che è stazionario, deterministico e di lunghezza incognita mediante la DFT, 
a partire dalla conoscenza di L campioni del segnale, pone i seguenti problemi: 

\begin{tcolorbox}[colback=red!5!white,colframe=red!75!black] \label{q: 45}
Come fa ad essere $x \left[ n\right]$ deterministico? Perchè la DFT viene fatta "off-line"? 
\end{tcolorbox}

\begin{itemize}
    \item risoluzione, cioè la distanza d in frequenza tra due sinusoidi, che dipende da L e dal tipo di finestra (che idealmente vorremmo avere uguale alla lunghezza di $x \left[ n\right]$, cioè infinita o molto lunga, ma che non è possibile perchè il tempo di osservazione per noi è L)
    \item leakage, la presenza dei lobi laterali che distorcono lo spettro (e quindi dipende solo dal tipo di finestra)
    \item granulosità, cioè l'effetto del campionamento in frequenza dovuto alla DFT; si può ridurre con zero padding, ma comporta problemi computazioni e dipende da N (e che dipende da quanti campioni abbiamo di $x \left[ n\right]$)
    \item precisione finita, dovuta al rumore nello spettro 
\end{itemize}

\newpage 

\section{Sequenze Random}

\footnote{ 
    Slide del prof | 13 - Stima spettrale | pag 11 - 13\\
    Video Lezione CADSP Lunedì 12 Dicembre 2022 | 1:32:50 - 1:42:20 
}

\begin{tcolorbox}
    Come abbiamo scritto precedentemnte, nella realtà non avremo a che fare con sequenze $x \left[ n\right]$ deterministiche, o di lunghezza finita o infinita, 
    bensì avremo a che fare con sequenze $x \left[ n\right]$ casuali: per questo, non prevedendo il futuro, l'unica cosa che possiamo fare a posteriori della sequenza casuale $x \left[ n\right]$ è studiarla con gli strumenti della statistica, 
    che abbiamo già studiato al corso di Teoria dei Segnali con il mitico prof. Franco Chiaraluce. \newline 

    Questa sezione è un ripasso di tutto quello che abbiamo studiato nei corsi precedenti. 
\end{tcolorbox}

Una sequenza casuale (o in inglese random) è una realizzazione di un processo stocastico del tipo: 

{
    \Large 
    \begin{equation}
        x \left[ n\right] = x_{n}
    \end{equation}
}

\begin{tcolorbox}
    Questa è la notazione utilizzata in questo corso di variabile statistica: nel corso di teoria di segnali, avremmo una notazione differente. 
\end{tcolorbox}

$x_{n}$ è una serie temporale; ad essa è generalmente associata una sua descrizione statistica del processo, per avere informazione su tutte le possibili realizzazione del processo. \newline 

\begin{tcolorbox}[colback=red!5!white,colframe=red!75!black] \label{q: 46}
$x_{n}$ è una serie temporale perchè varia nel tempo? 
\end{tcolorbox}

Tale descrizione statistica è completa quando conosciamo la distribuzione di probabilità del processo. \newline 

La descrizione diventa paraiale se conosciamo le medie di insieme del fenomeno $x_{n}$. \newline 

Di seguito i parametri che andremo ad utilizzare: 

\begin{itemize}
    \item media: 
    {
        \Large 
        \begin{equation}
            m_{x_{n}} = E \left[ x_{n}\right]
        \end{equation}
    }
    \item varianza: 
    {
        \Large 
        \begin{equation}
            \sigma_{x_{n}}^{2} = E \left[ \abs{x_{n} - m_{x_{n}}}\right]
        \end{equation}
    }
    \item autocorrelazione: 
    {
        \Large 
        \begin{equation}
            \sigma_{xx} \left[ n, m\right] = E \left[x_{n} \cdot x_{m}^{*}\right]
        \end{equation}
    } 
    \item autocovarianza: 
    {
        \Large 
        \begin{equation}
            \gamma_{xx} \left[ n, m\right] = E \left[ \left(x_{n} - m_{x_{n}}\right) \cdot \left(x_{m}^{*} - m_{x_{m}^{*}}\right)\right]
        \end{equation}
    }
\end{itemize}

\begin{tcolorbox}[colback=green!5!white,colframe=green!75!black]
TO-DO: Inserire dei video youtube o link per ripassare cosa significano i parametri statistici indicati. \newline 

Inoltre, inserire le formule degli integrali dal corso di Teoria dei segnali. 
\end{tcolorbox} 

\begin{tcolorbox}
Generalmente noi avremo a che fare con segnali a media nulla perchè sennò, in Fourier, avremo una Delta di Dirac in $n = 0$, che rappresenterebba la continua in Fourier.     
\end{tcolorbox}

Poichè, spesso non è possibile conoscere le media di insieme, ci si accontenta di conoscere le medie temporali della sequenza random a disposizione per processi stazionari in senso lato. \newline 

\begin{tcolorbox}[colback=green!5!white,colframe=green!75!black]
TO-DO: Aggiungere link per la definizione di processo stazionario e quale è la differenza in senso lato  
\end{tcolorbox} 

\begin{tcolorbox}[colback=red!5!white,colframe=red!75!black] \label{q: 47}
Differenza tra stazionarietà e stazionarietà in senso lato? 
\end{tcolorbox}

Allora, sicoccome nella realtà non conosciamo la PDF, non useremo i parametri appena indicati bensì i seguenti: 

\begin{itemize}
    \item media: 
    {
        \Large 
        \begin{equation}
            <x_{n}> = \lim_{L \to + \infty} \dfrac{1}{2 \cdot L + 1} \cdot \sum_{n = - L}^{L} x \left[ n\right] 
        \end{equation}
    }
\begin{tcolorbox}[colback=red!5!white,colframe=red!75!black] \label{q: 48}
Siccome L tende ad infinito, è come se la media di $x_{n}$ fosse calcolata fosse calcolata con una finestra $w \left[ n\right]$ infintia, giusto? 
\end{tcolorbox}
    \item vaianza: 
    {
        \Large 
        \begin{equation}
            <x_{n}>^{2} = \lim_{L \to + \infty} \dfrac{1}{2 \cdot L + 1} \cdot \sum_{n = - L}^{L} \abs{x \left[ n\right] - <x_{n}>}^{2} 
        \end{equation}
    }
    \item auto-correlazione: 
    {
        \Large 
        \begin{equation}
            <x_{n+m} \cdot x_{m}^{*}> = \phi_{xx} \left[ m\right] = \lim_{L \to + \infty} \dfrac{1}{2 \cdot L + 1} \cdot \sum_{n = - L}^{L} x_{n+m} \cdot x_{m}^{*}  
        \end{equation}
    }
\end{itemize}

Se le media temporali sono uguali a quelle di insieme, allora il processo si dice ergodico. \newline 

Inoltre, possiamo definire la cross-correlazione e la cross-covarianza tra due processi stocastici in maniera del tutto simile a quanto fatto per autocorrelazione ad autocovarianza 
(analogo discorso nel caso di medie temporali). \newline 

\begin{tcolorbox}[colback=red!5!white,colframe=red!75!black] \label{q: 49}
Cosa significa l'ultima frase? 
\end{tcolorbox}

Per una generica sequenza $x \left[ n\right]$, possiamo definire le seguenti quantità: 

\begin{itemize}
    \item energia di $x \left[ n\right]$: 
    {
        \Large 
        \begin{equation}
            e_{x} = \sum_{n = - \infty}^{+ \infty} \abs{x_{n}}^{2}
        \end{equation}
    }
    \item potenza di $x \left[ n\right]$: 
    {
        \Large 
        \begin{equation}
            p_{x} 
            = 
            \lim_{L \to + \infty}
            \dfrac{1}{2 \cdot L + 1} 
            \cdot 
            \sum_{n = - L}^{+ L}
            \abs{x_{n}}^{2}
        \end{equation}
    }
\end{itemize}

Se l'energia del segnale $x_{n}$ è finita, cioè: 

{
    \Large 
    \begin{equation}
        e_{x} < + \infty
    \end{equation}
}

allora, vale il Teorema di Parseval: 

{
    \Large 
    \begin{equation}
        e_{x} = \dfrac{1}{2 \pi} \cdot \int_{- \infty}^{+ \infty} \abs{X \left(e^{\jmath \omega}\right)}^{2} d \omega
    \end{equation}
}

dove: 

\begin{itemize}
    \item $X \left(e^{\jmath \omega}\right)$ è la trasformata di Fourier di $x_{n}$ 
    \item $\abs{X \left(e^{\jmath \omega}\right)}^{2}$ è lo spettro di densità di energia
\end{itemize}

Se i segnali $x_{n}$ sono ad energia finita, i segnali deterministici sono aperiodici; 
se sono a potenza finita, i segnali possono essere deterministici, periodici o random. \newline 

\begin{tcolorbox}[colback=red!5!white,colframe=red!75!black] \label{q: 50}
L'ultima frase: cosa? 
\end{tcolorbox}

\newpage 

\subsection{Spettro di Densità di potenza}

\footnote{ 
    Slide del prof | 13 - Stima spettrale | pag 15\\
    Appunti | 2025-12-16 | pag 1\\
    Video Lezione CADSP Lunedì 12 Dicembre 2022 | 1:42:20 - 1:48:40 
}

Per segnali a potenza finita (per cui non si ha la trasformata di Fourier perchè non è definibile direttamente), allora si fa uso delle definizioni di auto-correlazione a auto-covarianza, per cui la trasformata di Fourier esiste. \newline 

La trasformata di Fourier dell'autocorrelazione $\phi_{xx} \left[ m\right]$ è la seguente: 

{
    \Large 
    \begin{equation}
        P_{xx} \left(\omega\right) 
        = 
        \phi_{xx} \left(e^{\jmath \omega}\right)
        = 
        \sum_{n = - \infty}^{+ \infty}
        \phi_{xx} \left[ n\right]
        \cdot 
        e^{- \jmath \omega n}
        \quad 
        \text{ per }
        \abs{\omega} < \pi
    \end{equation}
}

che, $P_{xx} \left(\omega\right) $ va stimato. \newline 

Invece, la trasformata di Fouerier dell'autocovarianza $\phi_{xx} \left[ n, m\right]$ è la seguente: 

{
    \Large 
    \begin{equation}
        \Gamma_{xx} \left(e^{\jmath \omega}\right) 
        = 
        \sum_{n = - \infty}^{+ \infty} 
        \gamma_{xx} \left[ n\right] \cdot e^{- \jmath \omega n}
        \quad 
        \text{ per }
        \abs{\omega} < \pi 
    \end{equation}
}

Queste due formule di $P_{xx} \left(\omega\right)$ e $\Gamma_{xx} \left(e^{\jmath \omega}\right)$ sono molto utili perchè ci permettono di utilizzare strumenti deterministici 
per sequenze casuali, che non sono deterministici, senza la PDF. \newline 

In particolare, grazie al teorema di Wiener-Khintchine possiamo affermare che, per processi stazionari in senso lato ed ergodici, vale che: 

{
    \Large 
    \begin{equation}
        \text{Potenza di } x \left[ n\right] = p_{x} = \phi_{xx} \left[ 0\right] = \dfrac{1}{2 \pi} \cdot \int_{- \pi}^{+ \pi} P_{xx} \left(\omega\right) d\omega
    \end{equation}
}

Quindi, sotto queste ipotesi, il teorema di Wiener-Khintcine è equivalente al teorema di Parseval. \newline 

Spesso, l'ipotesi di ergodicità viene rilassata, cioè useremo il Teorema di Wiener-Khintichine anche per segnali non ergodici. \newline 

In modo analogo, si può definire lo spettro densità di potenza incrociata tra due segnali $x \left[ n\right]$ e $y \left[ n\right]$, come: 

{
    \Large 
    \begin{equation}
        P_{xy} \left(\omega\right)
        = 
        \Phi_{xy} \left(e^{\jmath \omega}\right) 
        = 
        \sum_{n = - \infty}^{+ \infty}
        \phi_{xy} \left[ n\right]
        \cdot 
        e^{- \jmath \omega }
        \quad 
        \text{ per }
        \abs{\omega} < \pi
    \end{equation}
}

\newpage 

\section{Stima dello spettro dei segnali random}

\footnote{ 
    Slide del prof | 13 - Stima spettrale | pag 16\\
    Appunti | 2025-12-16 | pag 1\\
    Video Lezione CADSP Lunedì 12 Dicembre 2022 | 1:54:35 - 2:00:55   
}

Il problema, come abbiamo sempre scritto, è quello di stimare lo spettro di potenza $P_{xx} \left(\omega\right)$ di un segnale discreto $x \left[ n\right]$ causale, ergodico e stazionario (in senso lato), 
a partire dalla conoscenza di un numero finito di valori di una realizzazione di $x \left[ n\right]$. \newline 

Poichè risulta che: 

{
    \Large 
    \begin{equation}
        P_{xx} \left(\omega\right)
        = 
        \sum_{n = - \infty}^{+ \infty}
        \phi_{xx} \left[ n\right]  \cdot e^{- \jmath \omega n}
    \end{equation}
}

che è una funzione reale sempre positiva, se $x \left[ n\right]$ è reale, allora lo spettro può essere calcolato solo conoscendo la sequenza di auto-correlazione $\phi_{xx} \left[ n\right]$, e quindi tutto il processo random. \newline 

La conoscenza di un numero limitato di campioni di realizzazione di $x \left[ n\right]$ permette solo di stimare lo spettro 
(lo stimiamo sempre perchè abbiamo solo campioni limitati di $x \left[ n\right]$ e non tutti i campioni per $n \in \left[ - \infty, + \infty\right]$). \newline 

Una buona stima di $P \left(\omega\right)$, cioè $\hat{P} \left(\omega\right)$ deve essere tale che: 

\begin{itemize}
    \item la stima non deve essere "polarizzata", cioè: 
    {
        \Large 
        \begin{equation}
            \text{valore atteso} 
            = 
            P_{xx} \left(\omega\right)
            \to 
            E \left[ \hat{P}_{xx} \left(\omega\right)\right] 
            = 
            P_{xx} \left(\omega\right)
        \end{equation}
    }
    \item la stima deve essere consistente: 
    {
        \Large 
        \begin{equation}
            \text{varianza }
            = 
            0 
            \to 
            E \left[ \abs{\hat{P}_{xx} \left(\omega\right) - E \left[ \hat{P}_{xx} \left(\omega\right)\right] }^{2}\right] 
            = 
            0
        \end{equation}
    } 
    cioè, che la variabile aleatoria $x \left[ n\right]$ sia il più possibile che deterministica, che abbia poca varianza, cioè vicina alla media, se il fenomeno è gaussiano, avremo una campana più stretta 
\end{itemize}

\begin{tcolorbox}[colback=red!5!white,colframe=red!75!black] \label{q: 51}
Ma queste due condizioni non le avremo mai perchè abbiamo a che fare con dei calcolatori, che usando degli LFSR, genereranno una sequenza polarizzata e non consistente? 
\end{tcolorbox} 

Queste relazioni dovrebbero valere per ogni pulsazione $\omega$ di $x \left[ n\right]$. \newline 

Le uguaglianze scritte, nella realtà, non saranno uguaglianze esatte, bensì i valori calcolati tenderenno o saranno simili alle formule. \newline 

All'aumentare del numero di campioni conosciuti di $x \left[ n\right]$, la stima dello spettro, cioè $\hat{P}_{xx} \left(\omega\right)$, deve tendere allo spettro vero, cioè $P_{xx} \left(\omega\right)$, con una variabilità sempre minore. \newline 

\newpage 

\section{Periodogramma}

\footnote{ 
    Slide del prof | 13 - Stima spettrale | pag 17 - 24\\
    Appunti | 2025-12-16 | pag 1\\
    Video Lezione CADSP Lunedì 12 Dicembre 2022 | 2:00:55  -  
}

Piuttosto che calcolarci la potenza $P_{xx} \left(\omega\right)$ attraverso l'autocorrelazione $\phi_{xx} \left[ n\right]$: 

{
    \Large 
    \begin{equation}
        P_{xx} \left(\omega\right)
        = 
        \sum_{n = - \infty}^{+ \infty}
        \phi_{xx} \left[ n\right]  \cdot e^{- \jmath \omega n}
    \end{equation}
}

l'idea è quella di stimare $P_{xx} \left(\omega\right)$ senza calcolarsi l'autocorrelazione $\phi_{xx} \left[ n\right]$. \newline 

